% ==STYLE-BEGIN==
\documentclass[UTF8,11pt]{ctexart}
\usepackage{amsmath,amssymb,amsthm}
\usepackage[a4paper,margin=2.4cm,top=2.3cm,bottom=2.3cm]{geometry}
\usepackage{fancyhdr}
\usepackage{enumitem}
\usepackage{hyperref}
\hypersetup{hidelinks}

% ---- 版式 ----
\linespread{1.15}
\setlength{\parskip}{0.35em}
\setlength{\parindent}{1.6em}
\setlist{nosep,leftmargin=2em}

% ---- 页眉 ----
\pagestyle{fancy}
\fancyhf{}
\fancyhead[L]{\small\coursename}
\fancyhead[R]{\small\thepage}
\renewcommand{\headrulewidth}{0.4pt}

% ---- 定理环境：中文用正体，标题加粗，按小节共用一条流水号 ----
\newtheoremstyle{cnthm}{0.7em}{0.7em}{\normalfont}{}{\bfseries}{}{0.5em}{}
\theoremstyle{cnthm}
\newtheorem{thm}{定理}[section]
\newtheorem{lem}[thm]{引理}
\newtheorem{prop}[thm]{命题}
\newtheorem{cor}[thm]{推论}
\newtheorem{defn}[thm]{定义}
\newtheorem{rem}[thm]{注}
\newtheorem{exa}[thm]{例}
\newtheorem{conj}[thm]{猜想}
\renewcommand{\proofname}{\bfseries 证明}
\providecommand{\coursename}{}
% ==STYLE-END==
\renewcommand{\coursename}{高等代数\,·\,1028}
\xeCJKDeclareCharClass{CJK}{"2150 -> "218F} % 让 Ⅰ Ⅱ 等罗马数字走中文字体
\title{高等代数 1028}
\date{}
\begin{document}
\maketitle

\section{线性方程组}

\begin{defn}[线性方程组]
形如
\[
\begin{cases}
a_{11}x_1 + a_{12}x_2 + \cdots + a_{1n}x_n = b_1,\\
a_{21}x_1 + a_{22}x_2 + \cdots + a_{2n}x_n = b_2,\\
\qquad\quad \cdots\cdots\\
a_{m1}x_1 + a_{m2}x_2 + \cdots + a_{mn}x_n = b_m
\end{cases}
\]
的方程组称为 $m$ 个方程 $n$ 个未知量的线性方程组，记作 (Ⅰ)。其中 $x_1,\cdots,x_n$ 是方程组的未知量，$a_{ij}$ 是第 $i$ 个方程中 $x_j$ 的系数，$b_1,\cdots,b_m$ 是常数。
\end{defn}

\begin{defn}[解、解集、同解]
如果 $x_1=k_1,\cdots,x_n=k_n$ 满足方程组 (Ⅰ)，则称 $x_1=k_1,\cdots,x_n=k_n$ 是方程组的解；全体解构成的集合称为解集。若方程组 (Ⅰ) 与 (Ⅱ) 的解集相同，则称 (Ⅰ) 与 (Ⅱ) 同解。
\end{defn}

\begin{rem}
解方程组要回答三个问题：（1）方程组是否有解；（2）若有解，如何求解；（3）解集的结构如何。

方程组可以有唯一解，可以无解，也可以有无穷多组解。
\end{rem}

\begin{defn}[初等变换]
（1）两个方程对换；（2）某个方程乘以非零常数；（3）某个方程的倍数加到另一个方程上去。以上三种变换称为方程组的初等变换。
\end{defn}

\begin{thm}[初等变换的性质]
（1）初等变换不改变方程组的解，即变换前后同解；
（2）初等变换是可逆的。
\end{thm}

\begin{defn}[阶梯形方程组]
称方程组为\textbf{阶梯形方程组}，是指
\begin{enumerate}
\item 不含未知量的等式永远在含未知量的等式的下方；
\item 含有未知量的方程中，第一个未知量（称为主未知量）的下标向下严格递增。
\end{enumerate}
\end{defn}

\begin{thm}[化阶梯形]
任何一个线性方程组都可以由初等变换化为阶梯形方程组。
\end{thm}

\begin{proof}
对方程的个数 $m$ 归纳。

（1）当 $m=1$ 时，显然是阶梯形方程组。

（2）假设方程个数为 $m-1$ 时命题成立。

（3）当方程个数为 $m$ 时，不妨设 $a_{11}\neq0$；否则存在 $a_{i1}\neq0$，此时把第 $1$ 个方程与第 $i$ 个方程对调。再把第一个方程的 $-\dfrac{a_{i1}}{a_{11}}$ 倍加到第 $i$ 个方程（$i\geq2$），得到
\[
\begin{cases}
a'_{11}x_1 + a'_{12}x_2 + \cdots + a'_{1n}x_n = b'_1,\\
\qquad a'_{22}x_2 + \cdots + a'_{2n}x_n = b'_2,\\
\qquad\qquad \cdots\cdots\\
\qquad\qquad\qquad a'_{m2}x_2 + \cdots + a'_{mn}x_n = b'_m.
\end{cases}
\]
其中后 $m-1$ 个方程（记为 $(*)$）只含 $x_2,\cdots,x_n$，共 $m-1$ 个方程。由归纳假设，$(*)$ 可化为阶梯形方程组；于是原方程组也可以化为阶梯形方程组。

因此，任何一个线性方程组都可由初等变换化为阶梯形方程组。
\end{proof}

\begin{rem}
上述过程就是 \textbf{Gauss 消元法}：用初等变换把方程组化为阶梯形，再逐层回代。
\end{rem}

\section{阶梯形方程组的类型}

\begin{enumerate}
\item 出现矛盾方程（$0=d\neq0$）：方程组无解。

\item 所有未知量都是主未知量：
\[
\begin{cases}
a_{11}x_1 + a_{12}x_2 + a_{13}x_3 + \cdots + a_{1n}x_n = b_1,\\
\qquad a_{22}x_2 + a_{23}x_3 + \cdots + a_{2n}x_n = b_2,\\
\qquad\qquad a_{33}x_3 + \cdots + a_{3n}x_n = b_3,\\
\qquad\qquad\qquad \ddots\\
\qquad\qquad\qquad\qquad a_{nn}x_n = b_n,\\
\qquad\qquad\qquad\qquad 0 = 0,\\
\qquad\qquad\qquad\qquad \vdots
\end{cases}
\]
从下往上分别解出 $x_n,x_{n-1},\cdots,x_1$，此时方程组有唯一解。

\item 设 $x_{i_1},x_{i_2},\cdots,x_{i_r}$ 是主未知量，剩余的未知量称为\textbf{自由未知量}，记为 $x_{j_1},\cdots,x_{j_{n-r}}$。此时方程组形如
\[
\begin{cases}
a'_{1i_1}x_{i_1} + a'_{1,i_1+1}x_{i_1+1} + \cdots + a'_{1n}x_n = b'_1,\\
\qquad a'_{2i_2}x_{i_2} + \cdots + a'_{2n}x_n = b'_2,\\
\qquad\qquad \cdots\cdots\\
\qquad\qquad\qquad a'_{ri_r}x_{i_r} + \cdots + a'_{rn}x_n = b'_r,\\
\qquad\qquad\qquad\qquad 0 = 0,\\
\qquad\qquad\qquad\qquad \vdots
\end{cases}
\]
把含自由未知量的项移到右边，得
\[
\begin{cases}
a'_{1i_1}x_{i_1} = b'_1 - a'_{1j_1}x_{j_1} - \cdots - a'_{1,j_{n-r}}x_{j_{n-r}},\\
a'_{2i_2}x_{i_2} = b'_2 - a'_{2j_1}x_{j_1} - \cdots - a'_{2,j_{n-r}}x_{j_{n-r}},\\
\qquad\cdots\cdots\\
a'_{ri_r}x_{i_r} = b'_r - a'_{rj_1}x_{j_1} - \cdots - a'_{r,j_{n-r}}x_{j_{n-r}}.
\end{cases}
\]
用自由未知量表示主未知量，即可解出方程组的全部解。
\end{enumerate}

\end{document}
